\chapter*{Введение}
\addcontentsline{toc}{chapter}{Введение}
\label{ch:intro}
Целью расчетно-графической работы является разработка клиент-серверного приложения, включающего Android-клиент для сбора геолокационных данных и параметров радиосигнала, а также backend-сервер на языке C++ для приёма, хранения в базе данных PostgreSQL, визуализации на OpenStreetMap картах и генерации тепловых карт на основе метода обратных взвешенных расстояний.

Для достижения поставленной цели в работе решаются следующие задачи:
\begin{itemize}
    \item Разработка Android-приложения, обеспечивающего получение данных о местоположении.
    \item Реализация в Android-приложении сбора информации о сотах сотовой связи стандартов LTE, GSM и NR посредством класса Telephony.
    \item Организация фонового сбора данных и передачи их на сервер через сокеты ZeroMQ.
    \item Разработка backend-сервера на языке C++ с многопоточной архитектурой: отдельный поток для работы ZMQ-сокетов и отдельный поток для графического интерфейса пользователя.
    \item Интеграция с СУБД PostgreSQL для сохранения всех входящих данных в структурированную таблицу.
    \item Реализация модуля загрузки и отображения тайлов OpenStreetMap с динамическим изменением в зависимости от положения камеры, включая кэширование загруженных изображений.
    \item Построение тепловых карт для критериев RSRP, RSRQ, RSSI с возможностью выбора отображаемого параметра через GUI.
\end{itemize}

В ходе работы были использованы следующие инструменты и технологии:
язык Kotlin для написания Android-приложения и C++ для серверной части;
библиотека ZeroMQ для обмена данными между телефоном и компьютером;
libpq для подключения к базе данных PostgreSQL;
libcurl для скачивания карт из интернета;
ImGui и ImPlot для создания окон и графиков;
GitHub для хранения кода и отслеживания изменений.

В результате получился работающий программно-аппаратный комплекс, который в реальном времени собирает данные о сигнале сотовой сети, помогает оценить качество покрытия и позволяет проводить замеры в городе с последующим анализом собранной информации.
\endinput